\section{Business context and problem}
\label{sec:context}

Cyclistic is a fictional bike-share company in Chicago whose data are those of the Divvy network: according to the case study, a fleet of 5,824 geotracked bicycles locked into a network of 692 stations, including assistive options (reclining bikes, hand tricycles, cargo bikes).
It offers single-ride passes, full-day passes and annual memberships.
According to the case study's finance analysts, \textbf{annual members are much more profitable} than casual riders.

Lily Moreno, the director of marketing, believes that future growth depends on maximizing the number of annual members and that there is an opportunity to convert casual riders, who are already aware of the service.
Three questions guide the future marketing program; this report addresses the first one.

\begin{keybox}[Assigned task]
\textbf{Question 1:} How do annual members and casual riders use Cyclistic bikes differently?\par
\smallskip
\textcolor{muted}{Questions 2 and 3, addressed by other analyses: Why would casual riders buy annual memberships? How can digital media influence them to become members?}
\end{keybox}

\textbf{Stakeholders.} Lily Moreno (project sponsor), the marketing analytics team, and the executive team, described as notoriously detail-oriented, which must approve the recommended program.

\textbf{Deliverables required by the case study.} Statement of the business task (Section~\ref{sec:context}); description of the data sources (Section~\ref{sec:data}); documentation of data cleaning (Section~\ref{sec:cleaning}); summary of the analysis (Section~\ref{sec:analysis}); visualizations and key findings (Section~\ref{sec:visuals}); three recommendations (Section~\ref{sec:recos}).
